\section{Calculo infinitesimal}
\subsection{Fórmulas con precisión de Orden 1: $\mathcal{O}(h)$}
El error de truncamiento depende linealmente del tamaño del paso $h$. 

\subsubsection{Operador en adelanto (Derivada primera)}
Utiliza el punto actual y el siguiente. Es la aproximación más básica.
\begin{equation}
    f'(x_0) = \frac{f(x_0+h) - f(x_0)}{h} - \underbrace{\frac{h}{2} \cdot f''(x_0)}_{\mathcal{O}(h)}
\end{equation}
\subsubsection{Operador en atraso (Derivada primera)}

$$f'(x_0) = \frac{f(x_0) - f(x_0-h)}{h} + \underbrace{\frac{h}{2} f''(x_0)}_{\mathcal{O}(h)}$$

\subsection{Fórmulas con precisión de Orden 2: $\mathcal{O}(h^2)$}
El error de truncamiento está elevado al cuadrado, por lo que si el paso $h$ se reduce a la mitad, el error se divide por cuatro.

\subsubsection{Operador de diferencias centrales (Derivada primera)}
Utiliza un punto hacia adelante y uno hacia atrás.
\begin{equation}
    f'(x_0) = \frac{f(x_0+h) - f(x_0-h)}{2h} - \underbrace{\frac{h^2}{3!} f'''(x_0)}_{\mathcal{O}(h^2)}
\end{equation}

\subsubsection{Operador de diferencias centrales (Derivada segunda)}
Utiliza el punto actual y sus dos vecinos inmediatos.
\begin{equation}
    f''(x_0) = \frac{f(x_0+h) - 2f(x_0) + f(x_0-h)}{h^2} - \underbrace{\frac{h^2}{12} f^{(4)}(x_0)}_{\mathcal{O}(h^2)}
\end{equation}

\subsubsection{Operador en adelanto de 3 puntos (Derivada primera)}
Ideal para calcular la derivada en el límite inferior de una tabla de datos donde no existen puntos previos (no se puede usar la central), manteniendo una alta precisión.
\begin{equation}
    f'(x_0) = \frac{-3 \cdot f(x_0) + 4 \cdot f(x_0+h) - f(x_0+2h)}{2 \cdot h} + \underbrace{\frac{2h^2}{3!} f'''(x_0)}_{\mathcal{O}(h^2)}
\end{equation}
Atraso
$$u_i' \approx \frac{3u_i - 4u_{i-1} + u_{i-2}}{2h}$$